\skyrius{Produkto patikra ir rezultatų vertinimas}
\subsection{Patikros metodika}
\pildyti{Kurią produkto versiją, kokioje aplinkoje ir kokiais duomenimis
tikrinote? Kaip atrinkote bandymus ar dalyvius? Kokį rezultatą iš anksto
laikėte tinkamu? Aprašymo turi pakakti bandymui suprasti ir, kai įmanoma, pakartoti.}
\gaires{Pasirinkite produktui tinkamus bandymus. Programai gali reikėti
įprastų, ribinių ir klaidingų įvesčių patikros; įrenginiui -- veikimo, tikslumo
ar patvarumo matavimų. Naudotojų atsiliepimų reikia tada, kai jie padeda
atsakyti į jūsų klausimą. Dalyvių ar bandymų skaičiaus šablonas nenustato.}

\subsection{Rezultatai pagal reikalavimus}
\begin{table}[H]
\caption{Produkto patikros rezultatų santrauka}
\label{tab:bandymai}
\begin{tabularx}{\textwidth}{@{}>{\raggedright\arraybackslash}p{2.2cm}YYY@{}}
\toprule
Bandymas / reikalavimas & Tikėtasi & Faktiškai gauta ir įrodymas & Įvertinimas \\
\midrule
T1 / R1 & \pildyti{Kriterijus} & \pildyti{Duomenys, įrašo ID} & \pildyti{Atitinka / ne} \\
T2 / R2 & \pildyti{Kriterijus} & \pildyti{Duomenys, įrašo ID} & \pildyti{Atitinka / ne} \\
\bottomrule
\end{tabularx}
\end{table}
\pildyti{Paaiškinkite \ref{tab:bandymai} lentelės rezultatus: ką jie leidžia
teigti apie produkto tinkamumą? Nesėkmingo bandymo atveju aprašykite priežastį,
pataisymą ir pakartotinės patikros rezultatą. Išsamūs įrašai -- \ref{app:testai} priede.}

\subsection{Nauda, ribotumai ir galutinė versija}
\pildyti{Kiek produktas išsprendžia įvade apibrėžtą problemą? Kokių išvadų
negalite daryti dėl bandymų sąlygų, imties ar neįgyvendintų funkcijų?
Palyginkite rezultatą su pradiniais reikalavimais ir aptartais analogais.}
\pildyti{Nurodykite galutinės versijos žymą ir pateikiamų failų vietą.
Naudojimo arba paleidimo instrukcija -- \ref{app:produktas} priede.
Vien veikiančios nuorodos nepakanka versijai apibrėžti: pridėkite vertinimui
skirtą kopiją arba kitą su vadovu suderintą produkto pateikimo būdą.}
